\onbaslik{ÖZET}
\begin{center}
{\setstretch{1.5}\bfseries\BaslikTR\par}
\end{center}
\vspace{0.4cm}

Üniversitelerde öğrenci destek birimlerine gelen soruların büyük bölümü, cevabı resmî yönetmelik
ve duyurularda zaten bulunan tekrar eden sorulardır; buna karşın hesap ele geçirilmesi, taciz,
not itirazı gibi hassas talepler mutlaka bir insan tarafından ele alınmalıdır. Bu tasarım
çalışmasının amacı, Bilecik Şeyh Edebali Üniversitesi öğrencilerinin Türkçe veya İngilizce yazdığı
destek taleplerini okuyan, kategori, öncelik ve insana yönlendirme gereği bakımından etiketleyen,
cevabı üniversitenin resmî belgelerinde arayan, kaynağını gösteren bir taslak cevap yazan ve bu
cevabın öğrenciye otomatik gönderilip gönderilmeyeceğine büyük dil modeli yerine test edilmiş, açık
kurallardan oluşan bir program koduyla karar veren bir yardım masası asistanı geliştirmektir.

Sistemde etiketleme ve taslak yazma adımlarında ücretsiz katmandaki bir büyük dil modeli, belge
aramada çok dilli metin gömmeleri ve vektör aramalı ilişkisel bir veritabanı kullanılmıştır. Yedi
koşulun tamamı sağlanmadıkça cevap gönderilmez; alıntılanan her kural, bulunan belgede kelimesi
kelimesine aranır. Sistem öğrenci sohbet ekranı, talep takibi ve birime göre yetkilendirilmiş
personel paneliyle birlikte canlıya alınmıştır.

Ara dönem itibarıyla elle etiketlenmiş yüz altmış talep üzerinde yapılan ölçümlerde, ayar için hiç
kullanılmayan iki kilitli test setinde insana gitmesi gereken taleplerden hiçbiri kaçırılmamış, doğru
belge ilk üç sonuç içinde her seferinde bulunmuş, gerçekçi sette hatalı otomatik gönderim
görülmemiş ve on iki saldırı denemesinin tamamı engellenmiştir. Üç etiketin birlikte doğruluğu
yaklaşık yüzde yetmiş üçtür. Kalan dönemde karma arama, daha büyük bir değerlendirme seti ve
modellerin karşılaştırılması planlanmaktadır.

\vspace{1cm}
\noindent{\singlespacing\textbf{Anahtar Kelimeler:} Büyük Dil Modelleri, Erişim Destekli Üretim,
Destek Talebi Sınıflandırma, Vektör Arama, Yardım Masası\par}

%%%%%%%%%%%%%%%%%%%%%%%%%%%%%%%%%%%%%%%%%%%%%%%%%%%%%%%%%%%%%%%%%%%%%%%%%%%%%%%%%
\onbaslik{ABSTRACT}
\begin{center}
{\setstretch{1.5}\bfseries\BaslikEN\par}
\end{center}
\vspace{0.4cm}

Most questions that reach university student support offices are repeated questions whose answers
already exist in official regulations and announcements, while sensitive requests such as a
compromised account, harassment or a grade appeal must always be handled by a person. The aim of
this design project is a helpdesk assistant for the students of Bilecik Şeyh Edebali University: it
reads a support request written in Turkish or English, labels its category, priority and whether a
person must take over, searches the official university documents for the answer, drafts a reply
that cites its sources, and decides whether the reply is sent automatically with tested, explicit
program code instead of the language model.

A free-tier large language model labels requests and writes drafts; multilingual text embeddings
and a relational database with vector search find the documents. A reply is sent only if all seven
conditions hold, and every quoted rule is searched for word for word in the retrieved document. The
system is deployed with a student chat page, request tracking and an office-scoped staff console.

At this interim stage, on one hundred and sixty hand-labelled requests, the two locked test sets
that were never used for tuning show no missed request that needed a person, the right document in
the top three results every time, no unsafe automatic reply on the realistic set, and all twelve
attack attempts blocked. All three labels are correct together for about seventy-three percent of
requests. Hybrid search, a larger evaluation set and a comparison of models are planned for the rest
of the term.

\vspace{1cm}
\noindent{\singlespacing\textbf{Keywords:} Large Language Models, Retrieval-Augmented Generation,
Support Ticket Classification, Vector Search, Helpdesk\par}
